\documentclass[11pt]{article}
\usepackage[a4paper,margin=1.9cm]{geometry}
\usepackage[T1]{fontenc}
\usepackage{textcomp}
\usepackage[spanish,es-noshorthands]{babel}
\usepackage{amsmath,amssymb}
\usepackage{enumitem}
\usepackage{tikz}
\usetikzlibrary{calc,arrows.meta,patterns,decorations.pathmorphing,shadows}
\usepackage{xcolor}
\usepackage{multicol}
\usepackage{parskip}
\usepackage{lmodern}
\renewcommand{\familydefault}{\sfdefault}

\definecolor{unalblue}{RGB}{31,73,124}
\definecolor{cajaazul}{RGB}{20,60,110}
\definecolor{cajaverde}{RGB}{35,120,75}
\definecolor{cajaroja}{RGB}{165,25,25}

\newlist{opciones}{enumerate}{1}
\setlist[opciones]{label=(\alph*),itemsep=1pt,topsep=2pt,leftmargin=1.8em,font=\normalfont}

\newcommand{\vect}[2]{\begin{pmatrix}#1\\#2\end{pmatrix}}
\newcommand{\vectiii}[3]{\begin{pmatrix}#1\\#2\\#3\end{pmatrix}}

\newcommand{\cajatitulada}[3]{%
  \begin{center}
  \begin{tikzpicture}
    \node[draw=#1, line width=2.2pt, rounded corners=4pt, inner sep=12pt,
          text width=0.92\linewidth, align=left,
          fill=#1!2.2]
      (box) {#3};
    \node[fill=white, anchor=west, xshift=10pt, inner sep=3pt] at (box.north west)
      {\small\color{#1}\bfseries #2};
  \end{tikzpicture}
  \end{center}
}

\newcommand{\examheader}{%
  \noindent
  \begin{minipage}[t]{0.6\linewidth}
    {\small\bfseries MÉTODOS NUMÉRICOS (3006907) --- Tercer Examen Parcial}\\
    {\small Semestre 01-2017, 5 de junio de 2017 --- TEMA A --- Con solución verificada}
  \end{minipage}%
  \hfill
  \begin{minipage}[t]{0.35\linewidth}
    \raggedleft\small Escuela de Matemáticas. Universidad Nacional de Colombia, Sede Medellín.
  \end{minipage}\\[2pt]
  \rule{\linewidth}{0.6pt}
}
\newcommand{\examfooter}{%
  \vfill
  \rule{\linewidth}{0.4pt}\\[1pt]
  {\scriptsize\textit{Transcripción digital --- MathUNAL}\hfill\textcolor{unalblue}{\textbf{\thepage}}}
}
\pagestyle{empty}

\begin{document}
\examheader

\vspace{4pt}
\noindent Nombre completo:\underline{\hspace{7.5cm}}\quad Carné:\underline{\hspace{2.5cm}}\quad Grupo:\underline{\hspace{1.5cm}}

\vspace{6pt}
\noindent\textit{La interpretación del examen hace parte de la evaluación, por tanto no se admiten preguntas comenzada la prueba.}

\vspace{6pt}
\noindent\fbox{\textbf{TODA RESPUESTA DEBE ESTAR DEBIDAMENTE JUSTIFICADA}}

\vspace{4pt}
\noindent Duración: 1 hora y 50 minutos.

\vspace{10pt}

\begin{enumerate}[label=\arabic*.,leftmargin=1.5em]

\item Considere el problema con valores iniciales (P.V.I.)
\[
\begin{cases}
y'(t) = \cos(ty(t)), \qquad -3 \le t \le 2.5, \\
y(-3) = -1.
\end{cases}
\]

\begin{enumerate}[label=\alph*.,leftmargin=1.8em]
\item (10\%) Demuestre que la función $f(t,y) := \cos(ty)$ satisface una condición de Lipschitz en la variable $y$ en la región $D := \{(t,y) \in \mathbb{R}^2 : -3 \le t \le 2.5, \ -\infty < y < \infty\}$.
\begin{itemize}[leftmargin=1.5em]
\item ¿Cuál es la constante de Lipschitz para esta función?
\item ¿Este P.V.I. tiene única solución?
\end{itemize}

\item (10\%) Para \textbf{este} P.V.I., escriba la fórmula en diferencias o fórmula de avance del método de Taylor 2.
\end{enumerate}

\cajatitulada{cajaazul}{Solución (a)}{
$f(t,y)$ es continua en $D$ (composición de funciones continuas), y $D$ es una región convexa (rectángulo infinito en $y$). Calculamos
\[
\frac{\partial f}{\partial y}(t,y) = -t\operatorname{sen}(ty).
\]
Como $|\operatorname{sen}(ty)| \le 1$ para todo $(t,y)$, se tiene
\[
\left|\frac{\partial f}{\partial y}(t,y)\right| = |t|\,|\operatorname{sen}(ty)| \le |t|.
\]
El máximo de $|t|$ en $[-3,2.5]$ se alcanza en $t=-3$, así que
\[
\boxed{L = 3}
\]
Como $f$ es continua en $D$, satisface condición de Lipschitz en $y$ con $L=3$, y $D$ es convexa, por el teorema de existencia y unicidad, \textbf{este P.V.I. tiene solución única}.
}

\vspace{6pt}
\cajatitulada{cajaazul}{Solución (b)}{
La fórmula de avance del método de Taylor de orden 2 es
\[
w_{i+1} = w_i + hf(t_i,w_i) + \frac{h^2}{2}f'(t_i,w_i), \qquad w_0 = \alpha,
\]
donde $f'(t,y)$ denota la \textbf{derivada total} de $f$ respecto a $t$ a lo largo de la solución, es decir
\[
f'(t,y) = \frac{\partial f}{\partial t} + \frac{\partial f}{\partial y}\cdot f(t,y).
\]
Calculamos cada parte:
\[
\frac{\partial f}{\partial t} = -y\operatorname{sen}(ty), \qquad \frac{\partial f}{\partial y} = -t\operatorname{sen}(ty)
\]
\[
f'(t,y) = -y\operatorname{sen}(ty) - t\operatorname{sen}(ty)\cdot\cos(ty) = -\operatorname{sen}(ty)\left(y + t\cos(ty)\right)
\]
Con $w_0=\alpha=-1$, la fórmula de avance queda:
\[
w_{i+1} = w_i + h\cos(t_iw_i) - \frac{h^2}{2}\operatorname{sen}(t_iw_i)\left(w_i + t_i\cos(t_iw_i)\right)
\]
}

\vspace{4pt}
\cajatitulada{cajaroja}{Nota de verificación}{
\small
En el manuscrito original, $f'(t,y)$ se calculó como $\dfrac{\partial f}{\partial t}+\dfrac{\partial f}{\partial y}$ (suma directa de las dos derivadas parciales), sin multiplicar el segundo término por $f(t,y)$. Para el método de Taylor de orden 2 aplicado a $y'=f(t,y)$, la fórmula requiere la \textbf{derivada total} $f'=\partial_t f + \partial_y f \cdot f$, ya que $y$ depende de $t$ a través de la propia ecuación diferencial. Arriba se dejó la expresión con ese factor incluido.
}

\vspace{10pt}

\item Considere el problema parabólico de la forma
\[
\begin{cases}
u_t(x,t) = \dfrac{1}{25}u_{xx}(x,t) + \dfrac{1}{5}u_x(x,t), & 0 < x < 1.2, \quad 0 < t < 0.2, \\[4pt]
u(x,0) = 1 - \left|x - \dfrac{3}{5}\right|, & 0 \le x \le 1.2, \\[4pt]
u(0,t) = \dfrac{2}{5} + t, & 0 \le t \le 0.2, \\[4pt]
u(1.2,t) = 3t + \dfrac{2}{5}, & 0 \le t \le 0.2.
\end{cases}
\]

Se quiere aproximar la solución empleando el método de diferencias finitas, con $h=\frac{1}{5}$, $k=0.04$.

\begin{enumerate}[label=\alph*.,leftmargin=1.8em]
\item (10\%) Discretizar el dominio (graficar la cuadrícula) e identificar en cuáles puntos $(x_i,t_j)$ se conoce la solución.

\item (10\%) Escribir la fórmula de diferencias finitas en la forma $w_{i,j+1} = \beta w_{i-1,j} + \gamma w_{i,j} + \eta w_{i+1,j}$, usando diferencia progresiva en $t$ y \textbf{centrada} en $x$.

\item (20\%) Hallar $A$, $B$, $C$ y $D$ en la tabla de aproximación.
\end{enumerate}

\cajatitulada{cajaazul}{Solución (b)}{
Partiendo de $\dfrac{w_{i,j+1}-w_{i,j}}{k} = \dfrac{1}{25}\cdot\dfrac{w_{i+1,j}-2w_{i,j}+w_{i-1,j}}{h^2} + \dfrac{1}{5}\cdot\dfrac{w_{i+1,j}-w_{i-1,j}}{2h}$ (diferencia \textbf{centrada} en $x$, tal como pide el enunciado), y despejando $w_{i,j+1}$ con $h=0.2$, $k=0.04$:
\[
w_{i,j+1} = \left(\frac{k}{25h^2}-\frac{k}{10h}\right)w_{i-1,j} + \left(1-\frac{2k}{25h^2}\right)w_{i,j} + \left(\frac{k}{25h^2}+\frac{k}{10h}\right)w_{i+1,j}
\]
Sustituyendo $h=0.2$, $k=0.04$:
\[
\boxed{w_{i,j+1} = 0.02\,w_{i-1,j} + 0.92\,w_{i,j} + 0.06\,w_{i+1,j}}
\]
válida para $i=1,\dots,N-1$ y $j=0,1,\dots,M-1$ (nodos interiores, no de frontera).
}

\vspace{4pt}
\cajatitulada{cajaroja}{Nota de verificación}{
\small
El manuscrito original usó diferencia \textbf{progresiva} para $u_x$ (es decir, $\frac{w_{i+1,j}-w_{i,j}}{h}$ en vez de $\frac{w_{i+1,j}-w_{i-1,j}}{2h}$), obteniendo $w_{i,j+1}=0.96\,w_{i,j}+0.04\,w_{i+1,j}$ (coeficiente $\beta=0$). El enunciado pide explícitamente diferencia \textbf{centrada} para las derivadas en $x$; con esa diferencia el coeficiente $\beta$ no se anula: la fórmula correcta lleva los tres coeficientes $0.02,\ 0.92,\ 0.06$ (arriba). Los incisos c) que siguen se resuelven con ambas fórmulas para que se pueda comparar contra el trabajo original del estudiante.
}

\vspace{6pt}
\cajatitulada{cajaverde}{Solución (c)}{
\textbf{Usando condiciones de frontera e inicial} (no dependen de la fórmula interior, así que $A$, $B$, $C$ coinciden en ambos enfoques):
\[
B = u(0.6,0) = 1-\left|0.6-\tfrac{3}{5}\right| = 1
\]
\[
A = u(0,t_3) = 0.4+0.12 = 0.52
\]
\[
C = u(1.2,t_4) = 3(0.16)+0.4 = 0.88
\]

\textbf{Para $D=w_{3,2}$} (nodo interior, sí depende de la fórmula usada), con $w_{3,1}=0.9840$:

Con la fórmula \textbf{centrada} (correcta, usando $w_{2,1}=0.8080$ y $w_{4,1}=0.7920$):
\[
D = 0.02(0.8080) + 0.92(0.9840) + 0.06(0.7920) \approx 0.9793
\]

Con la fórmula progresiva (la del manuscrito original, usando $w_{4,1}=0.7920$):
\[
D = 0.96(0.9840) + 0.04(0.7920) \approx 0.9763
\]

\[
\boxed{A=0.52 \qquad B=1 \qquad C=0.88 \qquad D\approx 0.98\ (\text{según fórmula usada})}
\]
}

\vspace{4pt}
\cajatitulada{cajaroja}{Nota de verificación}{
\small
En el manuscrito original, $B$ se calculó bien como $1$ pero luego se anotó en el recuadro final como \textbf{$B=-1$}, lo cual no coincide con el propio desarrollo (que da $1-|0.6-0.6|=1-0=1$). Para $D$, el manuscrito usó la recurrencia $w_{3,2}=0.96\,w_{3,1}-0.04\,w_{1,1}$: además del coeficiente ya señalado como no centrado, usa el vecino \textbf{izquierdo} $w_{1,1}$ con signo negativo en vez del vecino \textbf{derecho} $w_{4,1}$ con signo positivo — el resultado manuscrito ($0.91296$) no reproduce ninguna de las dos fórmulas correctamente.
}

\vspace{10pt}

\item (15\%) Considere el problema de contorno o de valores en la frontera (P.V.F.):
\[
\begin{cases}
y''(x) - 4x\,y'(x) + 3\operatorname{sen}(x)\,y(x) = \tan^{-1}(x), \qquad 2 < x < 3, \\
y(2) = -1.4, \\
y(3) = 2.4.
\end{cases}
\]

Enuncie los dos P.V.I.s que se deben resolver para utilizar el método del disparo, empleando $u$ y $v$. ¿Cómo se obtiene $y$ en términos de $u$ y $v$?

\cajatitulada{cajaazul}{Solución}{
\[
u:
\begin{cases}
u''(x) = 4x\,u'(x) - 3\operatorname{sen}(x)\,u(x) + \tan^{-1}(x), & 2 \le x \le 3, \\
u(2) = -1.4, \\
u'(2) = 0,
\end{cases}
\]
\[
v:
\begin{cases}
v''(x) = 4x\,v'(x) - 3\operatorname{sen}(x)\,v(x), & 2 \le x \le 3, \\
v(2) = 0, \\
v'(2) = 1.
\end{cases}
\]
Como $v(3) \ne 0$, la solución del P.V.F. se obtiene como
\[
y(x) = u(x) + \frac{2.4 - u(3)}{v(3)}\,v(x).
\]
}

\vspace{4pt}
\cajatitulada{cajaroja}{Nota de verificación}{
\small
El planteamiento del manuscrito original para $u$ y $v$ es correcto: despeja bien $y''$ de la ecuación original y arma los dos P.V.I. con las condiciones iniciales estándar del método del disparo ($u(2)=\alpha$, $u'(2)=0$; $v(2)=0$, $v'(2)=1$). No se detectaron errores en este punto.
}

\end{enumerate}

\examfooter
\end{document}
